\section*{अध्याय \ref{ch:g8:electric-current-effects} --- विद्युत धारा और उसके प्रभाव}

\begin{solution}{exo:g8:electric-current-effects:1}
ऊष्मीय (टोस्टर, केतली, तंतु); चुंबकीय (\omterm{def:g8:electric-current-effects:electromagnet}{विद्युत्चुंबक}, \omterm{ex:g6:circuits-and-safety:motor}{मोटर}, दरवाज़े
की घंटी); रासायनिक (चाँदी चढ़ाना, पानी को तोड़ना, बैटरियों का चार्ज
होना)।
\end{solution}

\begin{solution}{exo:g8:electric-current-effects:2}
ऊष्मीय; चुंबकीय; रासायनिक; ऊष्मीय; चुंबकीय।
\end{solution}

\begin{solution}{exo:g8:electric-current-effects:3}
यह कि धारा अपने तार के चारों ओर चुंबकत्व पैदा करती है --- यानी बिजली
और चुंबकत्व के बीच पहला पुल। धारा उलटने पर झटका भी उलट जाता है।
\end{solution}

\begin{solution}{exo:g8:electric-current-effects:4}
उसे बंद किया जा सकता है, चालू किया जा सकता है, और उसके ध्रुव मरज़ी से
पलटे जा सकते हैं --- और उसकी ताक़त की ख़ुराक भी तय की जा सकती है
(ज़्यादा फेरे, ज़्यादा धारा)। पत्थर का \omterm{def:g1:magnets:magnet}{चुंबक} इनमें से कुछ भी नहीं कर
सकता।
\end{solution}

\begin{solution}{exo:g8:electric-current-effects:5}
नंगे फेरे आपस में छूकर धारा को पूरी कुंडली घूमने के बजाय लपेटों के
बीच से ही शॉर्टकट दे देते; विद्युतरोधन उसे लंबा, \omterm{def:g1:magnets:magnet}{चुंबक} बनाने वाला रास्ता
लेने पर मजबूर करता है। और अकेली कुंडली उस क़तार का विरोध बहुत कम करती
है: जुड़ी छोड़ दो तो वह लगभग \omterm{def:g7:short-circuits-safety:device}{लघुपथन} है, जो कुंडली और \omterm{def:g3:first-electric-circuit:battery}{बैटरी} दोनों को
गरम कर देती है --- यहाँ \omterm{prop:g8:electric-current-effects:three}{ऊष्मीय प्रभाव} चौकसी करता है।
\end{solution}

\begin{solution}{exo:g8:electric-current-effects:6}
\omterm{ex:g3:first-electric-circuit:switch}{स्विच} खोलकर: धारा ग़ायब, चुंबकत्व ग़ायब, कार गिरी। उतनी ताक़त वाला
पत्थर का \omterm{def:g1:magnets:magnet}{चुंबक} तो हमेशा के लिए पकड़े रहता --- कार छुड़ाने के लिए उसे
खींचकर अलग करने वाली मशीनें माँगनी पड़तीं।
\end{solution}

\begin{solution}{exo:g8:electric-current-effects:7}
चलते हुए: संचित रासायनिक \omterm{def:g5:everyday-energy:energy}{ऊर्जा} विद्युत \omterm{def:g5:everyday-energy:energy}{ऊर्जा} बन जाती है, जो उस क़तार
को चलाती है (और आख़िर में \omterm{def:g1:light-and-shadows:source}{प्रकाश}, गति और गरमाहट बनकर ख़त्म होती है)।
चार्जर पर: विद्युत \omterm{def:g5:everyday-energy:energy}{ऊर्जा} सेल में से उलटी दिशा में मजबूर की जाकर
रासायनिक भंडार दोबारा खड़ा कर देती है।
\end{solution}

\begin{solution}{exo:g8:electric-current-effects:8}
क्योंकि ख़ुद वह क़तार अदृश्य है: उसका भेद सिर्फ़ उसके काम खोलते हैं।
गरम होना, झटका खाना और बुलबुले उठना ही धारा के अकेले दस्तख़त हैं ---
और अगले अध्याय का उपकरण चुंबकीय झटके को अंकों वाले पाठ में सुधार देता
है।
\end{solution}

\begin{solution}{exo:g8:electric-current-effects:9}
बटन दबा: धारा बही, \omterm{def:g8:electric-current-effects:electromagnet}{विद्युत्चुंबक} ने हथौड़ा झपटा, ठन --- पर हथौड़े की
वही \omterm{def:g7:motion-average-speed:formula}{चाल} उसकी अपनी आपूर्ति का एक संपर्क खोल देती है; चुंबकत्व मर जाता
है, कमानी हथौड़ा वापस लाती है, संपर्क फिर बन जाता है; धारा लौटती है,
फिर झपट --- और यही सिलसिला, हर सेकंड कई बार। पकड़े रहना डिज़ाइन से ही
नामुमकिन है: हर झपट अपनी ही वजह काट देती है --- और वह खड़खड़ाहट ही
घंटी की आवाज़ \emph{है}।
\end{solution}

\begin{solution}{exo:g8:electric-current-effects:10}
तापनियंत्रक के नाज़ुक संपर्क सिर्फ़ कुंडली की छोटी-सी धारा ढोते हैं;
और कुंडली का चुंबकत्व उन तगड़े संपर्कों को जोड़ता है जो एक अलग फंदे
में हीटर की भूखी धारा ढोते हैं। छोटा बड़े को हुक्म देता है, और दोनों
परिपथ कभी आपस में छूते नहीं --- नाज़ुक हाथ \omterm{def:g4:weight-and-mass:weight}{भार} नहीं, एक लीवर खींचता
है।
\end{solution}

\begin{solution}{exo:g8:electric-current-effects:11}
तरीक़ा: \omterm{def:g4:magnets-and-compass:compass}{दिक्सूचक} को दीवार से सटाकर सपाट पकड़ो, लोहे के फ़र्नीचर से
दूर; धीरे-धीरे घुमाते चलो; जहाँ शक़ वाला उपकरण चालू-बंद करने पर सुई
झटका खाए, वहीं धारा ढोता तार जा रहा है। तेज़ झटके का मतलब है पास वाली
या तगड़ी धारा। ईमानदार सीमा: सुइयाँ लोहे के पाइपों, चुंबकों और ख़ुद
पृथ्वी को भी जवाब देती हैं --- इसलिए धारा वही झटका साबित करता है जो
\omterm{ex:g3:first-electric-circuit:switch}{स्विच} के चालू-बंद के पीछे-पीछे चले।
\end{solution}

\begin{solution}{exo:g8:electric-current-effects:12}
धारा कुंडली को \omterm{def:g1:magnets:magnet}{चुंबक} बना देती है; कुंडली और ढाँचे के विपरीत ध्रुव
खिंचते हैं, समान ध्रुव धकेलते हैं, और कुंडली एक क़तार में आने की ओर
घूम जाती है। उलटाव के बिना वह उस क़तार पर पहुँचकर रुक जाती ---
आधा चक्कर, फिर काँपते हुए ठहराव: दरवाज़ा बंद करने वाला पुरज़ा, \omterm{ex:g6:circuits-and-safety:motor}{मोटर}
नहीं। क़तार वाले पल पर धारा उलट देना कुंडली के ध्रुव आपस में बदल देता
है: कल का खिंचाव अब धक्का है, और क़तार में आ जाना सबसे बुरी जगह बन
जाती है; इसलिए कुंडली को उस लक्ष्य का पीछा करते रहना पड़ता है जो हर
आधे चक्कर पर हटकर पलट जाता है --- यानी अनंत पीछा, जिसका मतलब है:
घूर्णन।
\end{solution}

\begin{solution}{pb:g8:electric-current-effects:1}
\textbf{1.} तीनों प्रदर्शन एक ही फंदे पर, \omterm{ex:g3:first-electric-circuit:switch}{स्विच} के साथ,
\emph{श्रेणी} में: तब एक ही क़तार बारी-बारी गरम तार, \omterm{def:g4:magnets-and-compass:compass}{दिक्सूचक} वाले
तार और मर्तबान की सेवा कर देती है --- एक \omterm{ex:g3:first-electric-circuit:switch}{स्विच}, तीन प्रभाव।
\textbf{2.} श्रेणी वाला नियम: मोटे और पतले, दोनों तारों को \emph{वही}
धारा पार करती है। फ़र्क़ सिर्फ़ सड़क में है --- पतला तार उस क़तार को
ज़्यादा झेलता है और गरम हो जाता है: ख़ुराक वही, मरीज़ अलग।
\textbf{3.} \omterm{def:g3:first-electric-circuit:battery}{बैटरी} उलटो \emph{और} तार \omterm{def:g4:magnets-and-compass:compass}{दिक्सूचक} की दूसरी ओर से ले जाओ
(या \omterm{def:g4:magnets-and-compass:compass}{दिक्सूचक} ही पलट दो): दोनों उलटाव आपस में कट जाते हैं, और झटका
पूरब की ओर ही बना रहता है।
\textbf{4.} यह कि पानी के दोनों घटक दो-और-एक के अनुपात में मौजूद हैं
--- एक सिरे से दूसरे के मुक़ाबले दोगुनी \omterm{def:g3:solids-liquids-gases:gas}{गैस}।
\textbf{5.} ख़ुद को रोकने वाला वह संपर्क हटाना पड़ेगा (वही पुरज़ा,
जिसे हथौड़े की \omterm{def:g7:motion-average-speed:formula}{चाल} खोल देती है): तब \omterm{def:g8:electric-current-effects:electromagnet}{विद्युत्चुंबक} फ़ेल्ट वाले हथौड़े
को ख़ामोशी से घंटी पर टिकाए रखेगा --- मेहमान को एक हल्की-सी ``दम''
सुनाई देगी और बस।
\textbf{6.} हौद पर चलाने वाला कुंडली का \omterm{ex:g3:first-electric-circuit:switch}{स्विच} खोल देता है: चुंबकत्व
ग़ायब, लोहा गिरा। एल्यूमिनियम आगे सवार निकल जाता है, क्योंकि \omterm{def:g1:magnets:magnet}{चुंबक} ---
चाहे पत्थर के हों या धारा से बने --- सिर्फ़ लोहे के कुनबे को पकड़ते
हैं, और एल्यूमिनियम ने उन्हें कभी जवाब दिया ही नहीं।
\textbf{7.} यह कि धारा से बने \omterm{def:g1:magnets:magnet}{चुंबक} के ध्रुव धारा की दिशा के पीछे
चलते हैं: एक उलटो, दूसरा उलट जाता है।
\textbf{8.} \omterm{def:g3:first-electric-circuit:battery}{बैटरी} चाँदी चढ़ाने का काम अपना रासायनिक भंडार
\emph{ख़र्च} करके चलाती है; और किसी भंडार को दोबारा खड़ा करने के लिए
बाहर से मजबूर की गई \omterm{def:g5:everyday-energy:energy}{ऊर्जा} चाहिए। जो \omterm{def:g3:first-electric-circuit:battery}{बैटरी} उसी काम से ख़ुद को भर ले
जिसका दाम वह चुका रही है, वह फिर वही हमेशा चलने वाली मशीन है ---
\omterm{def:g5:everyday-energy:energy}{ऊर्जा} का बिल ख़ुद अपना दाम नहीं चुका सकता।
\textbf{9.} पकड़: \omterm{ex:g3:first-electric-circuit:switch}{स्विच} बंद, कुंडली \omterm{def:g1:magnets:magnet}{चुंबक} बनी, चीज़ पकड़ में। झूलना:
धारा बनी रही, पकड़ बनी रही --- \omterm{prop:g8:electric-current-effects:three}{चुंबकीय प्रभाव} ड्यूटी पर। छोड़:
नाली के ऊपर \omterm{ex:g3:first-electric-circuit:switch}{स्विच} खुला, चुंबकत्व मरा, इनाम गिरा।
\textbf{10.} यह कि \omterm{def:g8:electric-current-effects:electromagnet}{विद्युत्चुंबक} अपनी धारा के दम पर जीता है: \omterm{def:g3:first-electric-circuit:battery}{बैटरी}
ढीली, क़तार कमज़ोर, पकड़ कमज़ोर --- चालू-बंद होने की सुविधा का दाम है
यही निर्भरता। (पत्थर का \omterm{def:g1:magnets:magnet}{चुंबक} कभी ढीला नहीं पड़ता --- और कभी छोड़ता
भी नहीं।)
\textbf{11.} उतनी ही धारा पर ज़्यादा फेरे, ज़्यादा मज़बूत \omterm{def:g1:magnets:magnet}{चुंबक}: यानी
पक्की पकड़। दूसरी ख़ुराक ख़ुद धारा की ताक़त है --- यानी \emph{तीव्रता},
जिसे अगले अध्याय में मापा जाएगा।
\textbf{12.} जैसे: ``ऊष्मीय: \omterm{def:g7:short-circuits-safety:fuse}{फ़्यूज़}, जो सबसे पहले और वफ़ादारी से
मरता है। चुंबकीय: क्रेन-खेल, जो हुक्म पर पकड़ता और छोड़ता है।
रासायनिक: चाँदी चढ़ाने वाला मर्तबान, जो \omterm{def:g3:first-electric-circuit:battery}{बैटरी} के सब्र से चम्मचों को
चाँदी पहना देता है।''
\end{solution}
